% Part: many-valued-logic
% Chapter: infinite-valued-logics
% Section: introduction

\documentclass[../../../include/open-logic-section]{subfiles}

\begin{document}

\olfileid[fr]{mvl}{inf}{int}

\olsection{Introduction}

Le nombre de valeurs de vérité d'une
matrice n'est pas nécessairement fini.
Un choix naturel pour un ensemble
infini de valeurs de vérité est celui
des nombres rationnels compris entre
$0$ et~$1$,
$V_\infty=[0,1]\cap\Rat$;
autrement dit,
\begin{align*}
  V_\infty
    &=\Setabs{\frac{n}{m}}{n,m\in\Nat,
      \ m\ge1,\ 0\le n\le m}.
\intertext{Lorsque nous considérons
cet ensemble infini de valeurs de vérité,
il est souvent utile d'étudier aussi,
pour tout entier $m\ge2$, les sous-ensembles}
  V_m
    &=\Setabs{\frac{n}{m-1}}{n\in\Nat,
      \ 0\le n\le m-1}.
\intertext{Les $m$ valeurs de $V_m$ sont régulièrement espacées
de $0$ à $1$. Pour $m=5$, on obtient :}
  V_5
    &=\{0,\frac14,\frac12,
      \frac34,1\}.
\end{align*}
Dans les logiques fondées sur ces
ensembles de valeurs, seule la
valeur~$1$ est généralement désignée,
c'est-à-dire $V^+=\{1\}$.
Autrement dit, $1$ joue le rôle
de la vérité absolue et $0$ celui
de la fausseté absolue, mais les
formules peuvent prendre toute valeur
intermédiaire de~$V$.

Nous pouvons aussi considérer
$V_{[0,1]}=[0,1]$, l'ensemble de
tous les nombres \emph{réels} entre
$0$ et~$1$, ou d'autres sous-ensembles
infinis de~$[0,1]$. Les logiques
ayant $V_{[0,1]}$ pour ensemble de
valeurs de vérité sont souvent
appelées \emph{logiques floues}.

\end{document}
